\QHeading{第8题　哪些是轴对称图形？在（\hspace{25pt}）里画“$\checkmark$”。}

\par\noindent\makebox[\linewidth][c]{
\begin{tikzpicture}[x=1bp,y=-1bp,color=black,draw=black,text=black]
\begin{scope}\clip (56.0,120.67) rectangle (515.0,269.33);
\hdhead{99}{综合测试}
\hdsection{8}{148}{哪些是轴对称图形？在（\hspace{25pt}）里画“$\checkmark$”。}
\draw[line width=.65pt] (124.000,144.000) ellipse [x radius=17.7bp,y radius=13bp];
\begin{scope}[scale=.6666667]
\draw[line width=.7pt] (315,209) .. controls (298,204) and (320,191) .. (325,195)
 .. controls (333,182) and (343,190) .. (334,198)
 .. controls (347,199) and (343,207) .. (333,209)
 .. controls (361,214) and (353,236) .. (328,236)
 .. controls (299,237) and (296,218) .. (315,209);
\draw[line width=.7pt] (430,211) rectangle (499,235);
\draw[line width=.7pt] (569,214) .. controls (562,209) and (577,190) .. (596,191)
 .. controls (616,192) and (628,201) .. (630,212)
 .. controls (632,224) and (612,211) .. (603,209)
 .. controls (604,225) and (599,235) .. (586,236)
 .. controls (578,235) and (583,229) .. (589,228)
 .. controls (597,224) and (596,216) .. (596,210)
 .. controls (583,216) and (575,220) .. (569,214);
\begin{scope}[shift={(732,207)}]
\draw[line width=.65pt] (0,-17)..controls(-7,-32)and(-11,-20)..(-8,-18)
 ..controls(-20,-23)and(-16,-11)..(-9,-10)
 ..controls(-21,-9)and(-12,1)..(-5,-1)
 ..controls(-10,11)and(-1,11)..(0,2)
 ..controls(1,11)and(10,11)..(5,-1)
 ..controls(12,1)and(21,-9)..(9,-10)
 ..controls(16,-11)and(20,-23)..(8,-18)
 ..controls(11,-20)and(7,-32)..(0,-17);
\draw (0,-17)--(0,28);
\draw (0,16)..controls(-20,2)and(-24,15)..(0,20);
\draw (0,16)..controls(20,2)and(24,15)..(0,20);
\end{scope}
\begin{scope}[shift={(185,302)}]
\draw (0,-14) ellipse (13 and 10);
\draw (-8,-4)..controls(-26,4)and(-26,21)..(-8,18)
 --(-8,30)..controls(-8,37)and(-1,37)..(0,31)
 ..controls(1,37)and(8,37)..(8,30)--(8,18)
 ..controls(26,21)and(26,4)..(8,-4);
\draw (-7,-17)..controls(-5,-19)and(-3,-19)..(-2,-17)
 (2,-17)..controls(3,-19)and(5,-19)..(7,-17);
\draw (-6,-11)..controls(-2,-7)and(2,-7)..(6,-11);
\end{scope}
\begin{scope}[shift={(326,319)}]
\path[fill=black] (-25,-2)..controls(-23,-16)and(-8,-22)..(-1,-20)
 --(4,-24)..controls(16,-23)and(25,-16)..(28,-10)
 --(23,-7)--(7,-20)--(4,-19)--(21,-4)--(15,2)--(28,2)
 ..controls(25,17)and(10,24)..(-6,20)
 ..controls(-22,17)and(-29,8)..(-25,-2);
\path[fill=white] (-12,0)--(0,-10)--(12,0)--(0,10)--cycle;
\path[fill=white] (2,-19)--(5,-21)--(24,-6)--(22,-4)--cycle;
\end{scope}
\foreach \x/\y in {465/298,445/314,485/314,465/330}{
\draw (\x,\y-8)--(\x+10,\y)--(\x,\y+8)--(\x-10,\y)--cycle;}
\draw (456,306)--(473,321) (473,307)--(456,322);
\draw (598,288)--(612,300)--(605,300)--(619,313)--(610,313)
 --(628,326)--(604,326)--(604,339)--(592,339)--(592,326)
 --(570,326)--(586,313)--(578,313)--(591,300)--(584,300)--cycle;
\begin{scope}[shift={(733,315)}]
\foreach \a in {0,90,180,270}{\begin{scope}[rotate=\a]
\draw (0,0)--(0,-24)..controls(18,-24)and(18,-6)..(0,0);
\end{scope}}\end{scope}
\end{scope}
\q{151}{249}{105}{12}{（\hspace{28pt}）}
\q{293}{249}{105}{12}{（\hspace{28pt}）}
\q{429}{249}{105}{12}{（\hspace{28pt}）}
\q{565}{249}{105}{12}{（\hspace{28pt}）}
\q{698}{249}{105}{12}{（\hspace{28pt}）}
\q{151}{351}{105}{12}{（\hspace{28pt}）}
\q{293}{351}{105}{12}{（\hspace{28pt}）}
\q{429}{351}{105}{12}{（\hspace{28pt}）}
\q{565}{351}{105}{12}{（\hspace{28pt}）}
\q{698}{351}{105}{12}{（\hspace{28pt}）}
\hdsection{9}{407}{动手操作。}
\q{95}{451}{440}{12}{① 过点$A$画已知直线的垂线。}
\q{95}{494}{440}{12}{② 过点$A$画已知直线的平行线。}
\fill (436.000,316.667) circle (1.2bp);
\node[inner sep=0pt,align=center,font=\normalsize,] at (434.000,303.333) {$A$};
\draw[line width=.7pt] (434.667,349.333)--(498.000,332.000);
\hdsection{10}{570}{下面是四（1）班第一、二小组同学的一次数学测试成绩。（单}
\q{127}{609}{700}{13.5}{\color{bookred}位：分）}
\q{96}{642}{124}{12}{第一小组：}
\q{490}{642}{124}{12}{第二小组：}
\node[inner sep=0pt,align=center,font=\normalsize,] at (141.333,434.667) {84};
\node[inner sep=0pt,align=center,font=\normalsize,] at (174.667,434.667) {75};
\node[inner sep=0pt,align=center,font=\normalsize,] at (208.000,434.667) {85};
\node[inner sep=0pt,align=center,font=\normalsize,] at (240.667,434.667) {70};
\node[inner sep=0pt,align=center,font=\normalsize,] at (274.000,434.667) {100};
\node[inner sep=0pt,align=center,font=\normalsize,] at (404.000,434.667) {92};
\node[inner sep=0pt,align=center,font=\normalsize,] at (436.667,434.667) {95};
\node[inner sep=0pt,align=center,font=\normalsize,] at (470.000,434.667) {87};
\node[inner sep=0pt,align=center,font=\normalsize,] at (502.667,434.667) {72};
\node[inner sep=0pt,align=center,font=\normalsize,] at (536.000,434.667) {83};
\node[inner sep=0pt,align=center,font=\normalsize,] at (141.333,460.000) {62};
\node[inner sep=0pt,align=center,font=\normalsize,] at (174.667,460.000) {85};
\node[inner sep=0pt,align=center,font=\normalsize,] at (208.000,460.000) {98};
\node[inner sep=0pt,align=center,font=\normalsize,] at (240.667,460.000) {94};
\node[inner sep=0pt,align=center,font=\normalsize,] at (274.000,460.000) {88};
\node[inner sep=0pt,align=center,font=\normalsize,] at (404.000,460.000) {60};
\node[inner sep=0pt,align=center,font=\normalsize,] at (436.667,460.000) {100};
\node[inner sep=0pt,align=center,font=\normalsize,] at (470.000,460.000) {95};
\node[inner sep=0pt,align=center,font=\normalsize,] at (502.667,460.000) {91};
\node[inner sep=0pt,align=center,font=\normalsize,] at (536.000,460.000) {80};
\node[inner sep=0pt,align=center,font=\normalsize,] at (141.333,485.333) {86};
\node[inner sep=0pt,align=center,font=\normalsize,] at (174.667,485.333) {55};
\node[inner sep=0pt,align=center,font=\normalsize,] at (208.000,485.333) {70};
\node[inner sep=0pt,align=center,font=\normalsize,] at (240.667,485.333) {81};
\node[inner sep=0pt,align=center,font=\normalsize,] at (274.000,485.333) {97};
\node[inner sep=0pt,align=center,font=\normalsize,] at (404.000,485.333) {90};
\node[inner sep=0pt,align=center,font=\normalsize,] at (436.667,485.333) {72};
\node[inner sep=0pt,align=center,font=\normalsize,] at (470.000,485.333) {87};
\node[inner sep=0pt,align=center,font=\normalsize,] at (502.667,485.333) {93};
\node[inner sep=0pt,align=center,font=\normalsize,] at (536.000,485.333) {93};
\q{95}{756}{737}{12}{① 分段整理上面的数据，然后填入下表。}
\q{254}{795}{570}{11.7}{四（1）班第一、二小组同学数学测试成绩统计表}
\path[fill=bookpink!70] (84.667,551.333) rectangle (156.667,580.667);
\node[inner sep=0pt,align=center,font=\normalsize,text width=68.667bp] at (120.667,566.000) {成绩/分};
\path[fill=bookpink!70] (156.667,551.333) rectangle (206.667,580.667);
\node[inner sep=0pt,align=center,font=\normalsize,text width=46.667bp] at (181.667,566.000) {合计};
\path[fill=bookpink!70] (206.667,551.333) rectangle (266.667,580.667);
\node[inner sep=0pt,align=center,font=\normalsize,text width=56.667bp] at (236.667,566.000) {60以下};
\path[fill=bookpink!70] (266.667,551.333) rectangle (326.000,580.667);
\node[inner sep=0pt,align=center,font=\normalsize,text width=56.000bp] at (296.333,566.000) {60～69};
\path[fill=bookpink!70] (326.000,551.333) rectangle (384.667,580.667);
\node[inner sep=0pt,align=center,font=\normalsize,text width=55.333bp] at (355.333,566.000) {70～79};
\path[fill=bookpink!70] (384.667,551.333) rectangle (444.000,580.667);
\node[inner sep=0pt,align=center,font=\normalsize,text width=56.000bp] at (414.333,566.000) {80～89};
\path[fill=bookpink!70] (444.000,551.333) rectangle (501.333,580.667);
\node[inner sep=0pt,align=center,font=\normalsize,text width=54.000bp] at (472.667,566.000) {90～99};
\path[fill=bookpink!70] (501.333,551.333) rectangle (549.333,580.667);
\node[inner sep=0pt,align=center,font=\normalsize,text width=44.667bp] at (525.333,566.000) {100};
\node[inner sep=0pt,align=center,font=\normalsize,text width=68.667bp] at (120.667,595.000) {总计};
\node[inner sep=0pt,align=center,font=\normalsize,text width=68.667bp] at (120.667,623.667) {第一小组};
\node[inner sep=0pt,align=center,font=\normalsize,text width=68.667bp] at (120.667,652.667) {第二小组};
\draw[line width=.5pt] (84.667,551.333)--(84.667,667.333);
\draw[line width=.5pt] (156.667,551.333)--(156.667,667.333);
\draw[line width=.5pt] (206.667,551.333)--(206.667,667.333);
\draw[line width=.5pt] (266.667,551.333)--(266.667,667.333);
\draw[line width=.5pt] (326.000,551.333)--(326.000,667.333);
\draw[line width=.5pt] (384.667,551.333)--(384.667,667.333);
\draw[line width=.5pt] (444.000,551.333)--(444.000,667.333);
\draw[line width=.5pt] (501.333,551.333)--(501.333,667.333);
\draw[line width=.5pt] (549.333,551.333)--(549.333,667.333);
\draw[line width=.5pt] (84.667,551.333)--(549.333,551.333);
\draw[line width=.5pt] (84.667,580.667)--(549.333,580.667);
\draw[line width=.5pt] (84.667,609.333)--(549.333,609.333);
\draw[line width=.5pt] (84.667,638.000)--(549.333,638.000);
\draw[line width=.5pt] (84.667,667.333)--(549.333,667.333);
\q{95}{1018}{390}{12}{② 第一小组的平均成绩是多少？}
\q{126}{1060}{345}{12}{第二小组呢？}
\q{488}{1018}{349}{12}{③ 两个小组的平均成绩是多少？}

\begin{scope}[color=answerred,draw=answerred,text=answerred]\ans{118.667}{165.333}{30.000}{13}{$\checkmark$}
\ans{305.333}{165.333}{30.000}{13}{$\checkmark$}
\ans{483.333}{165.333}{30.000}{13}{$\checkmark$}
\ans{118.667}{233.333}{30.000}{13}{$\checkmark$}
\ans{305.333}{233.333}{30.000}{13}{$\checkmark$}
\ans{394.000}{233.333}{30.000}{13}{$\checkmark$}
\begin{scope}[scale=0.6666667]\draw[answerred,line width=.85pt](618,484.853)--(753,447.906);\draw[answerred,line width=.85pt](654,475)--(674.210,548.850);\node[answerred,above]at(730,451){②};\node[answerred,right]at(673,543){①};\coordinate(foot)at(666.47,520.59);\draw[answerred](664.10,511.93)--(672.76,509.56)--(675.13,518.22);\end{scope}
\ans{178.000}{590.667}{37.333}{11}{$30$}
\ans{230.000}{590.667}{37.333}{11}{$1$}
\ans{288.000}{590.667}{37.333}{11}{$2$}
\ans{344.667}{590.667}{37.333}{11}{$5$}
\ans{404.000}{590.667}{37.333}{11}{$10$}
\ans{464.667}{590.667}{37.333}{11}{$10$}
\ans{518.667}{590.667}{37.333}{11}{$2$}
\ans{178.000}{619.333}{37.333}{11}{$15$}
\ans{230.000}{619.333}{37.333}{11}{$1$}
\ans{288.000}{619.333}{37.333}{11}{$1$}
\ans{344.667}{619.333}{37.333}{11}{$3$}
\ans{404.000}{619.333}{37.333}{11}{$6$}
\ans{464.667}{619.333}{37.333}{11}{$3$}
\ans{518.667}{619.333}{37.333}{11}{$1$}
\ans{178.000}{648.000}{37.333}{11}{$15$}
\ans{230.000}{648.000}{37.333}{11}{$0$}
\ans{288.000}{648.000}{37.333}{11}{$1$}
\ans{344.667}{648.000}{37.333}{11}{$2$}
\ans{404.000}{648.000}{37.333}{11}{$4$}
\ans{464.667}{648.000}{37.333}{11}{$7$}
\ans{518.667}{648.000}{37.333}{11}{$1$}
\ans{84.000}{719}{244}{9.5}{第一小组：$1230\div15=82$（分）。\\答：第一小组的平均成绩是82分。\\第二小组：$1290\div15=86$（分）。\\答：第二小组的平均成绩是86分。}
\ans{341.333}{715.333}{223}{9.5}{{\setlength{\jot}{-2pt}$\begin{aligned}[t]&\phantom{={}}(1230+1290)\div30\\&=2520\div30\\&=84\text{（分）}\end{aligned}$}\\答：两个小组的平均成绩是84分。}
\ans{84.000}{776}{460.000}{9.5}{本题两组总分的逐项计算，见紧接本页的计算续页。}
\end{scope}
\end{scope}\pgfresetboundingbox\path[use as bounding box] (56.0,120.67) rectangle (515.0,269.33);
\end{tikzpicture}}\par
\par\noindent




\par

\par\noindent

\QHeading{第9题　动手操作。}

\par\noindent\makebox[\linewidth][c]{
\begin{tikzpicture}[x=1bp,y=-1bp,color=black,draw=black,text=black]
\begin{scope}\clip (56.0,293.33) rectangle (507.0,378.0);
\hdhead{99}{综合测试}
\hdsection{8}{148}{哪些是轴对称图形？在（\hspace{25pt}）里画“$\checkmark$”。}
\draw[line width=.65pt] (124.000,144.000) ellipse [x radius=17.7bp,y radius=13bp];
\begin{scope}[scale=.6666667]
\draw[line width=.7pt] (315,209) .. controls (298,204) and (320,191) .. (325,195)
 .. controls (333,182) and (343,190) .. (334,198)
 .. controls (347,199) and (343,207) .. (333,209)
 .. controls (361,214) and (353,236) .. (328,236)
 .. controls (299,237) and (296,218) .. (315,209);
\draw[line width=.7pt] (430,211) rectangle (499,235);
\draw[line width=.7pt] (569,214) .. controls (562,209) and (577,190) .. (596,191)
 .. controls (616,192) and (628,201) .. (630,212)
 .. controls (632,224) and (612,211) .. (603,209)
 .. controls (604,225) and (599,235) .. (586,236)
 .. controls (578,235) and (583,229) .. (589,228)
 .. controls (597,224) and (596,216) .. (596,210)
 .. controls (583,216) and (575,220) .. (569,214);
\begin{scope}[shift={(732,207)}]
\draw[line width=.65pt] (0,-17)..controls(-7,-32)and(-11,-20)..(-8,-18)
 ..controls(-20,-23)and(-16,-11)..(-9,-10)
 ..controls(-21,-9)and(-12,1)..(-5,-1)
 ..controls(-10,11)and(-1,11)..(0,2)
 ..controls(1,11)and(10,11)..(5,-1)
 ..controls(12,1)and(21,-9)..(9,-10)
 ..controls(16,-11)and(20,-23)..(8,-18)
 ..controls(11,-20)and(7,-32)..(0,-17);
\draw (0,-17)--(0,28);
\draw (0,16)..controls(-20,2)and(-24,15)..(0,20);
\draw (0,16)..controls(20,2)and(24,15)..(0,20);
\end{scope}
\begin{scope}[shift={(185,302)}]
\draw (0,-14) ellipse (13 and 10);
\draw (-8,-4)..controls(-26,4)and(-26,21)..(-8,18)
 --(-8,30)..controls(-8,37)and(-1,37)..(0,31)
 ..controls(1,37)and(8,37)..(8,30)--(8,18)
 ..controls(26,21)and(26,4)..(8,-4);
\draw (-7,-17)..controls(-5,-19)and(-3,-19)..(-2,-17)
 (2,-17)..controls(3,-19)and(5,-19)..(7,-17);
\draw (-6,-11)..controls(-2,-7)and(2,-7)..(6,-11);
\end{scope}
\begin{scope}[shift={(326,319)}]
\path[fill=black] (-25,-2)..controls(-23,-16)and(-8,-22)..(-1,-20)
 --(4,-24)..controls(16,-23)and(25,-16)..(28,-10)
 --(23,-7)--(7,-20)--(4,-19)--(21,-4)--(15,2)--(28,2)
 ..controls(25,17)and(10,24)..(-6,20)
 ..controls(-22,17)and(-29,8)..(-25,-2);
\path[fill=white] (-12,0)--(0,-10)--(12,0)--(0,10)--cycle;
\path[fill=white] (2,-19)--(5,-21)--(24,-6)--(22,-4)--cycle;
\end{scope}
\foreach \x/\y in {465/298,445/314,485/314,465/330}{
\draw (\x,\y-8)--(\x+10,\y)--(\x,\y+8)--(\x-10,\y)--cycle;}
\draw (456,306)--(473,321) (473,307)--(456,322);
\draw (598,288)--(612,300)--(605,300)--(619,313)--(610,313)
 --(628,326)--(604,326)--(604,339)--(592,339)--(592,326)
 --(570,326)--(586,313)--(578,313)--(591,300)--(584,300)--cycle;
\begin{scope}[shift={(733,315)}]
\foreach \a in {0,90,180,270}{\begin{scope}[rotate=\a]
\draw (0,0)--(0,-24)..controls(18,-24)and(18,-6)..(0,0);
\end{scope}}\end{scope}
\end{scope}
\q{151}{249}{105}{12}{（\hspace{28pt}）}
\q{293}{249}{105}{12}{（\hspace{28pt}）}
\q{429}{249}{105}{12}{（\hspace{28pt}）}
\q{565}{249}{105}{12}{（\hspace{28pt}）}
\q{698}{249}{105}{12}{（\hspace{28pt}）}
\q{151}{351}{105}{12}{（\hspace{28pt}）}
\q{293}{351}{105}{12}{（\hspace{28pt}）}
\q{429}{351}{105}{12}{（\hspace{28pt}）}
\q{565}{351}{105}{12}{（\hspace{28pt}）}
\q{698}{351}{105}{12}{（\hspace{28pt}）}
\hdsection{9}{407}{动手操作。}
\q{95}{451}{440}{12}{① 过点$A$画已知直线的垂线。}
\q{95}{494}{440}{12}{② 过点$A$画已知直线的平行线。}
\fill (436.000,316.667) circle (1.2bp);
\node[inner sep=0pt,align=center,font=\normalsize,] at (434.000,303.333) {$A$};
\draw[line width=.7pt] (434.667,349.333)--(498.000,332.000);
\hdsection{10}{570}{下面是四（1）班第一、二小组同学的一次数学测试成绩。（单}
\q{127}{609}{700}{13.5}{\color{bookred}位：分）}
\q{96}{642}{124}{12}{第一小组：}
\q{490}{642}{124}{12}{第二小组：}
\node[inner sep=0pt,align=center,font=\normalsize,] at (141.333,434.667) {84};
\node[inner sep=0pt,align=center,font=\normalsize,] at (174.667,434.667) {75};
\node[inner sep=0pt,align=center,font=\normalsize,] at (208.000,434.667) {85};
\node[inner sep=0pt,align=center,font=\normalsize,] at (240.667,434.667) {70};
\node[inner sep=0pt,align=center,font=\normalsize,] at (274.000,434.667) {100};
\node[inner sep=0pt,align=center,font=\normalsize,] at (404.000,434.667) {92};
\node[inner sep=0pt,align=center,font=\normalsize,] at (436.667,434.667) {95};
\node[inner sep=0pt,align=center,font=\normalsize,] at (470.000,434.667) {87};
\node[inner sep=0pt,align=center,font=\normalsize,] at (502.667,434.667) {72};
\node[inner sep=0pt,align=center,font=\normalsize,] at (536.000,434.667) {83};
\node[inner sep=0pt,align=center,font=\normalsize,] at (141.333,460.000) {62};
\node[inner sep=0pt,align=center,font=\normalsize,] at (174.667,460.000) {85};
\node[inner sep=0pt,align=center,font=\normalsize,] at (208.000,460.000) {98};
\node[inner sep=0pt,align=center,font=\normalsize,] at (240.667,460.000) {94};
\node[inner sep=0pt,align=center,font=\normalsize,] at (274.000,460.000) {88};
\node[inner sep=0pt,align=center,font=\normalsize,] at (404.000,460.000) {60};
\node[inner sep=0pt,align=center,font=\normalsize,] at (436.667,460.000) {100};
\node[inner sep=0pt,align=center,font=\normalsize,] at (470.000,460.000) {95};
\node[inner sep=0pt,align=center,font=\normalsize,] at (502.667,460.000) {91};
\node[inner sep=0pt,align=center,font=\normalsize,] at (536.000,460.000) {80};
\node[inner sep=0pt,align=center,font=\normalsize,] at (141.333,485.333) {86};
\node[inner sep=0pt,align=center,font=\normalsize,] at (174.667,485.333) {55};
\node[inner sep=0pt,align=center,font=\normalsize,] at (208.000,485.333) {70};
\node[inner sep=0pt,align=center,font=\normalsize,] at (240.667,485.333) {81};
\node[inner sep=0pt,align=center,font=\normalsize,] at (274.000,485.333) {97};
\node[inner sep=0pt,align=center,font=\normalsize,] at (404.000,485.333) {90};
\node[inner sep=0pt,align=center,font=\normalsize,] at (436.667,485.333) {72};
\node[inner sep=0pt,align=center,font=\normalsize,] at (470.000,485.333) {87};
\node[inner sep=0pt,align=center,font=\normalsize,] at (502.667,485.333) {93};
\node[inner sep=0pt,align=center,font=\normalsize,] at (536.000,485.333) {93};
\q{95}{756}{737}{12}{① 分段整理上面的数据，然后填入下表。}
\q{254}{795}{570}{11.7}{四（1）班第一、二小组同学数学测试成绩统计表}
\path[fill=bookpink!70] (84.667,551.333) rectangle (156.667,580.667);
\node[inner sep=0pt,align=center,font=\normalsize,text width=68.667bp] at (120.667,566.000) {成绩/分};
\path[fill=bookpink!70] (156.667,551.333) rectangle (206.667,580.667);
\node[inner sep=0pt,align=center,font=\normalsize,text width=46.667bp] at (181.667,566.000) {合计};
\path[fill=bookpink!70] (206.667,551.333) rectangle (266.667,580.667);
\node[inner sep=0pt,align=center,font=\normalsize,text width=56.667bp] at (236.667,566.000) {60以下};
\path[fill=bookpink!70] (266.667,551.333) rectangle (326.000,580.667);
\node[inner sep=0pt,align=center,font=\normalsize,text width=56.000bp] at (296.333,566.000) {60～69};
\path[fill=bookpink!70] (326.000,551.333) rectangle (384.667,580.667);
\node[inner sep=0pt,align=center,font=\normalsize,text width=55.333bp] at (355.333,566.000) {70～79};
\path[fill=bookpink!70] (384.667,551.333) rectangle (444.000,580.667);
\node[inner sep=0pt,align=center,font=\normalsize,text width=56.000bp] at (414.333,566.000) {80～89};
\path[fill=bookpink!70] (444.000,551.333) rectangle (501.333,580.667);
\node[inner sep=0pt,align=center,font=\normalsize,text width=54.000bp] at (472.667,566.000) {90～99};
\path[fill=bookpink!70] (501.333,551.333) rectangle (549.333,580.667);
\node[inner sep=0pt,align=center,font=\normalsize,text width=44.667bp] at (525.333,566.000) {100};
\node[inner sep=0pt,align=center,font=\normalsize,text width=68.667bp] at (120.667,595.000) {总计};
\node[inner sep=0pt,align=center,font=\normalsize,text width=68.667bp] at (120.667,623.667) {第一小组};
\node[inner sep=0pt,align=center,font=\normalsize,text width=68.667bp] at (120.667,652.667) {第二小组};
\draw[line width=.5pt] (84.667,551.333)--(84.667,667.333);
\draw[line width=.5pt] (156.667,551.333)--(156.667,667.333);
\draw[line width=.5pt] (206.667,551.333)--(206.667,667.333);
\draw[line width=.5pt] (266.667,551.333)--(266.667,667.333);
\draw[line width=.5pt] (326.000,551.333)--(326.000,667.333);
\draw[line width=.5pt] (384.667,551.333)--(384.667,667.333);
\draw[line width=.5pt] (444.000,551.333)--(444.000,667.333);
\draw[line width=.5pt] (501.333,551.333)--(501.333,667.333);
\draw[line width=.5pt] (549.333,551.333)--(549.333,667.333);
\draw[line width=.5pt] (84.667,551.333)--(549.333,551.333);
\draw[line width=.5pt] (84.667,580.667)--(549.333,580.667);
\draw[line width=.5pt] (84.667,609.333)--(549.333,609.333);
\draw[line width=.5pt] (84.667,638.000)--(549.333,638.000);
\draw[line width=.5pt] (84.667,667.333)--(549.333,667.333);
\q{95}{1018}{390}{12}{② 第一小组的平均成绩是多少？}
\q{126}{1060}{345}{12}{第二小组呢？}
\q{488}{1018}{349}{12}{③ 两个小组的平均成绩是多少？}

\begin{scope}[color=answerred,draw=answerred,text=answerred]\ans{118.667}{165.333}{30.000}{13}{$\checkmark$}
\ans{305.333}{165.333}{30.000}{13}{$\checkmark$}
\ans{483.333}{165.333}{30.000}{13}{$\checkmark$}
\ans{118.667}{233.333}{30.000}{13}{$\checkmark$}
\ans{305.333}{233.333}{30.000}{13}{$\checkmark$}
\ans{394.000}{233.333}{30.000}{13}{$\checkmark$}
\begin{scope}[scale=0.6666667]\draw[answerred,line width=.85pt](618,484.853)--(753,447.906);\draw[answerred,line width=.85pt](654,475)--(674.210,548.850);\node[answerred,above]at(730,451){②};\node[answerred,right]at(673,543){①};\coordinate(foot)at(666.47,520.59);\draw[answerred](664.10,511.93)--(672.76,509.56)--(675.13,518.22);\end{scope}
\ans{178.000}{590.667}{37.333}{11}{$30$}
\ans{230.000}{590.667}{37.333}{11}{$1$}
\ans{288.000}{590.667}{37.333}{11}{$2$}
\ans{344.667}{590.667}{37.333}{11}{$5$}
\ans{404.000}{590.667}{37.333}{11}{$10$}
\ans{464.667}{590.667}{37.333}{11}{$10$}
\ans{518.667}{590.667}{37.333}{11}{$2$}
\ans{178.000}{619.333}{37.333}{11}{$15$}
\ans{230.000}{619.333}{37.333}{11}{$1$}
\ans{288.000}{619.333}{37.333}{11}{$1$}
\ans{344.667}{619.333}{37.333}{11}{$3$}
\ans{404.000}{619.333}{37.333}{11}{$6$}
\ans{464.667}{619.333}{37.333}{11}{$3$}
\ans{518.667}{619.333}{37.333}{11}{$1$}
\ans{178.000}{648.000}{37.333}{11}{$15$}
\ans{230.000}{648.000}{37.333}{11}{$0$}
\ans{288.000}{648.000}{37.333}{11}{$1$}
\ans{344.667}{648.000}{37.333}{11}{$2$}
\ans{404.000}{648.000}{37.333}{11}{$4$}
\ans{464.667}{648.000}{37.333}{11}{$7$}
\ans{518.667}{648.000}{37.333}{11}{$1$}
\ans{84.000}{719}{244}{9.5}{第一小组：$1230\div15=82$（分）。\\答：第一小组的平均成绩是82分。\\第二小组：$1290\div15=86$（分）。\\答：第二小组的平均成绩是86分。}
\ans{341.333}{715.333}{223}{9.5}{{\setlength{\jot}{-2pt}$\begin{aligned}[t]&\phantom{={}}(1230+1290)\div30\\&=2520\div30\\&=84\text{（分）}\end{aligned}$}\\答：两个小组的平均成绩是84分。}
\ans{84.000}{776}{460.000}{9.5}{本题两组总分的逐项计算，见紧接本页的计算续页。}
\end{scope}
\end{scope}\pgfresetboundingbox\path[use as bounding box] (56.0,293.33) rectangle (507.0,378.0);
\end{tikzpicture}}\par



\par



\par\noindent

\QHeading{第8、9题　补充检查方法}

\Answer{轴对称图形可用对折法或对应点法检查。过A作已知直线的垂线，可用三角板直接作，也可折叠描回；作平行线，可平移三角板，也可借助两次作垂线。过指定点向同一条直线作垂线或平行线，所得直线各只有一条。}

\QHeading{第10题　下面是四（1）班第一、二小组同学的一次数学测试成绩。（单位：分）}

\Problem{第一小组：}
\par

\keygrid{5}{84 & 75 & 85 & 70 & 100\\
62 & 85 & 98 & 94 & 88\\
86 & 55 & 70 & 81 & 97\\}
\Problem{第二小组：}
\keygrid{5}{92 & 95 & 87 & 72 & 83\\
60 & 100 & 95 & 91 & 80\\
90 & 72 & 87 & 93 & 93\\}
\Problem{① 分段整理上面的数据，然后填入下表。四（1）班第一、二小组同学数学测试成绩统计表}
\DataTable{4}{成绩/分 & 总计 & 第一小组 & 第二小组\\\hline
合计 & {\color{answerred}$30$} & {\color{answerred}$15$} & {\color{answerred}$15$}\\\hline
60以下 & {\color{answerred}$1$} & {\color{answerred}$1$} & {\color{answerred}$0$}\\\hline
60～69 & {\color{answerred}$2$} & {\color{answerred}$1$} & {\color{answerred}$1$}\\\hline
70～79 & {\color{answerred}$5$} & {\color{answerred}$3$} & {\color{answerred}$2$}\\\hline
80～89 & {\color{answerred}$10$} & {\color{answerred}$6$} & {\color{answerred}$4$}\\\hline
90～99 & {\color{answerred}$10$} & {\color{answerred}$3$} & {\color{answerred}$7$}\\\hline
100 & {\color{answerred}$2$} & {\color{answerred}$1$} & {\color{answerred}$1$}\\\hline}
\Problem{② 第一小组的平均成绩是多少？第二小组呢？}
\Answer{第一小组：$1230\div15=82$（分）。\par 答：第一小组的平均成绩是82分。\par 第二小组：$1290\div15=86$（分）。\par 答：第二小组的平均成绩是86分。}
\Problem{③ 两个小组的平均成绩是多少？}
\Answer{{\Formula{\begin{aligned}[t]&\phantom{={}}(1230+1290)\div30\\&=2520\div30\\&=84\text{（分）}\end{aligned}}}\par 答：两个小组的平均成绩是84分。}




\QHeading{多出的总分}

\Answer{对应高于80分的同学；等于80分的多出量为0。分组相加只是改变合并顺序，各组不重复、不遗漏。}

\QHeading{第1组}

\Answer{
\Formula{\begin{aligned}[t]&\phantom{={}}4+5+20\\
&=9+20\\
&=29\text{（分）}\end{aligned}}}

\QHeading{第2组}

\Answer{
\Formula{\begin{aligned}[t]&\phantom{={}}5+18+14\\
&=23+14\\
&=37\text{（分）}\end{aligned}}}

\QHeading{第3组}

\Answer{
\Formula{\begin{aligned}[t]&\phantom{={}}8+6+1\\
&=14+1\\
&=15\text{（分）}\end{aligned}}}

\QHeading{第4组}

\Answer{
\Formula{\begin{aligned}[t]&\phantom{={}}17\text{（分）}\end{aligned}}}

\QHeading{各组合计}

\Answer{
\Formula{\begin{aligned}[t]&\phantom{={}}29+37+15+17\\
&=66+15+17\\
&=81+17\\
&=98\text{（分）}\end{aligned}}}

\QHeading{不足的总分}

\Answer{对应低于80分的同学。分组相加只是改变合并顺序，各组不重复、不遗漏。}

\QHeading{第1组}

\Answer{
\Formula{\begin{aligned}[t]&\phantom{={}}5+10+18\\
&=15+18\\
&=33\text{（分）}\end{aligned}}}

\QHeading{第2组}

\Answer{
\Formula{\begin{aligned}[t]&\phantom{={}}25+10\\
&=35\text{（分）}\end{aligned}}}

\QHeading{各组合计}

\Answer{
\Formula{\begin{aligned}[t]&\phantom{={}}33+35\\
&=68\text{（分）}\end{aligned}}}

\QHeading{求平均成绩}

\Answer{把相抵后的多出量平均分给15人，再加80分。
\Formula{\begin{aligned}[t]&\phantom{={}}80+(98-68)\div 15\\
&=80+30\div 15\\
&=80+2\\
&=82\text{（分）}\end{aligned}}\par 答：第1小组平均成绩是82分。}

\QHeading{多出的总分}

\Answer{对应高于80分的同学；等于80分的多出量为0。分组相加只是改变合并顺序，各组不重复、不遗漏。}

\QHeading{第1组}

\Answer{
\Formula{\begin{aligned}[t]&\phantom{={}}12+15+7\\
&=27+7\\
&=34\text{（分）}\end{aligned}}}

\QHeading{第2组}

\Answer{
\Formula{\begin{aligned}[t]&\phantom{={}}3+20+15\\
&=23+15\\
&=38\text{（分）}\end{aligned}}}

\QHeading{第3组}

\Answer{
\Formula{\begin{aligned}[t]&\phantom{={}}11+0+10\\
&=11+10\\
&=21\text{（分）}\end{aligned}}}

\QHeading{第4组}

\Answer{
\Formula{\begin{aligned}[t]&\phantom{={}}7+13+13\\
&=20+13\\
&=33\text{（分）}\end{aligned}}}

\QHeading{各组合计}

\Answer{
\Formula{\begin{aligned}[t]&\phantom{={}}34+38+21+33\\
&=72+21+33\\
&=93+33\\
&=126\text{（分）}\end{aligned}}}

\QHeading{不足的总分}

\Answer{对应低于80分的同学。
\Formula{\begin{aligned}[t]&\phantom{={}}8+20+8\\
&=28+8\\
&=36\text{（分）}\end{aligned}}}

\QHeading{求平均成绩}

\Answer{把相抵后的多出量平均分给15人，再加80分。
\Formula{\begin{aligned}[t]&\phantom{={}}80+(126-36)\div 15\\
&=80+90\div 15\\
&=80+6\\
&=86\text{（分）}\end{aligned}}\par 答：第2小组平均成绩是86分。}

\QHeading{第10题②　补充方法：第1小组以80分为基数}

\Answer{每组都是15人。先把每人成绩看成80分，多出的分数与不足的分数分别相加。}

\QHeading{第10题②　补充方法：第2小组以80分为基数}

\Answer{每组都是15人。先把每人成绩看成80分，多出的分数与不足的分数分别相加。}

\QHeading{第10题③　补充方法：平均两个组的平均分}

\Answer{本题两个小组人数相同，都是15人，因此可以把82分和86分相加再除以2。人数不同时不能照用这个办法。
\Formula{\begin{aligned}[t]&\phantom{={}}(82+86)\div 2\\
&=168\div 2\\
&=84\text{（分）}\end{aligned}}\par 答：两个小组的平均成绩是84分。}

\QHeading{第10题①　另一种计数方法}

\Answer{除逐个画“正”字外，也可以先排序、再按分数段圈出连续的一组数。第一小组排序为55、62、70、70、75、81、84、85、85、86、88、94、97、98、100；第二小组排序为60、72、72、80、83、87、87、90、91、92、93、93、95、95、100。按表头逐段计数，所得各格与前表一致。}

\Answer{
本页列出原书第99页第10题两组总分的每次加法及平均数计算。
\QHeading{计算过程：第一小组总分与平均分}
\Formula{
\begin{aligned}[t]&\phantom{={}}84+75+85+70+100+62+85+98+94+88+86+55+70+81+97\\&=159+85+70+100+62+85+98+94+88+86+55+70+81+97\\&=244+70+100+62+85+98+94+88+86+55+70+81+97\\&=314+100+62+85+98+94+88+86+55+70+81+97\\&=414+62+85+98+94+88+86+55+70+81+97\\&=476+85+98+94+88+86+55+70+81+97\\&=561+98+94+88+86+55+70+81+97\\&=659+94+88+86+55+70+81+97\\&=753+88+86+55+70+81+97\\&=841+86+55+70+81+97\\&=927+55+70+81+97\\&=982+70+81+97\\&=1052+81+97\\&=1133+97\\&=1230\text{（分）}\end{aligned}
}
\Formula{
\begin{aligned}[t]&\phantom{={}}1230\div15\\&=82\text{（分）}\end{aligned}
}
答：第一小组的平均成绩是82分。

\QHeading{计算过程：第二小组总分与平均分}
\Formula{
\begin{aligned}[t]&\phantom{={}}92+95+87+72+83+60+100+95+91+80+90+72+87+93+93\\&=187+87+72+83+60+100+95+91+80+90+72+87+93+93\\&=274+72+83+60+100+95+91+80+90+72+87+93+93\\&=346+83+60+100+95+91+80+90+72+87+93+93\\&=429+60+100+95+91+80+90+72+87+93+93\\&=489+100+95+91+80+90+72+87+93+93\\&=589+95+91+80+90+72+87+93+93\\&=684+91+80+90+72+87+93+93\\&=775+80+90+72+87+93+93\\&=855+90+72+87+93+93\\&=945+72+87+93+93\\&=1017+87+93+93\\&=1104+93+93\\&=1197+93\\&=1290\text{（分）}\end{aligned}
}
\Formula{
\begin{aligned}[t]&\phantom{={}}1290\div15\\&=86\text{（分）}\end{aligned}
}
答：第二小组的平均成绩是86分。
\QHeading{计算过程：两组同学的平均成绩}
\Formula{
\begin{aligned}[t]&\phantom{={}}(1230+1290)\div(15+15)\\&=2520\div30\\&=84\text{（分）}\end{aligned}
}
答：两个小组同学的平均成绩是84分。

}
